\documentclass[11pt]{article}
\usepackage[margin=1in]{geometry}
\usepackage{amsmath,amssymb,amsthm,mathtools,bm}
\usepackage{microtype}
\usepackage[T1]{fontenc}
\usepackage{lmodern}
\usepackage{booktabs,enumitem}
\usepackage[hidelinks]{hyperref}
\setlist{nosep}
\newtheorem{theorem}{Theorem}[section]
\newtheorem{corollary}[theorem]{Corollary}
\newtheorem{proposition}[theorem]{Proposition}
\theoremstyle{definition}
\newtheorem{definition}[theorem]{Definition}
\newtheorem{assumption}[theorem]{Assumption}
\theoremstyle{remark}
\newtheorem{remark}[theorem]{Remark}

\newcommand{\Mb}{M_{\mathrm{b}}}
\newcommand{\vf}{v_{\mathrm{f}}}

\title{\textbf{Synopsis: BTFR Zero Point from Cosmological Horizon Memory}\\[0.35em]
\large What the theorems say (without the proofs)}
\author{Andrew Bruce Boyd}
\date{September 2026}

\begin{document}
\maketitle
\noindent\textit{Computational note.}
Proofs, exploratory experiments, and paper writing were assisted by generative AI.
Mathematical theorems stated in this paper are formally verified in Lean~4 in the
companion specifications, as faithful ports of the TeX arguments at the abstractions
recorded there, except results explicitly tagged as \emph{literature hypotheses}
(standard theorems cited from the literature and not re-proved here) or
\emph{physical inputs} (modeling assumptions outside mathematics).

\begin{abstract}
Short guide to the zero-point paper: galactic modelling already gives the BTFR
\emph{exponent}; cosmological-horizon Hysterical memory fixes
$a_0=H_0c/2\pi$.
Proofs live in the full paper.
\end{abstract}

\section{Why a zero point is needed}
Galactic modelling produces $\vf^{4}=K\Mb$ on a capacity attractor, with
$K=C^{2}A^{2}/y_{\ast}$.
Writing $\vf^{4}=Ga_0\Mb$ names $a_0=K/G$ but does not compute it.
This paper supplies that absolute scale from residual-gravity memory cut off by the
cosmological horizon.

\section{Horizon memory clock}
Euclidean regularity at a Killing horizon gives period $\beta=2\pi/\kappa$.
For the cosmological horizon, $\kappa=H_0$.

\begin{definition}[IR memory acceleration]
\[
  a_{\mathrm{IR}}:=\frac{c}{\beta}=\frac{c\kappa}{2\pi}.
\]
\end{definition}

\begin{theorem}[Cosmological memory scale]
\[
  \boxed{a_{\mathrm{IR}}=\frac{H_0c}{2\pi}.}
\]
\end{theorem}

Raw Hubble time would have given $cH_0$; the $2\pi$ is horizon KMS regularity.

\section{Subject independence}
\begin{proposition}
$a_{\mathrm{IR}}$ is a property of the residual channel and the FLRW congruence.
It does not depend on any particular probe's age or a passive basis rename.
\end{proposition}

\section{Pinning and matching}
Far-halo flat curves need a tiny loss rate; the softest allowed gravitational gap in
expansion is the cosmological horizon memory rate (physical assumption).
On the mature attractor one identifies $a_0=a_{\mathrm{IR}}$ (physical assumption).

\begin{theorem}[BTFR zero point]
\[
  \boxed{a_0=\frac{H_0c}{2\pi},}
  \qquad
  \vf^{4}=G\frac{H_0c}{2\pi}\,\Mb.
\]
\end{theorem}

\begin{corollary}[Numerics]
At Planck-like $H_0$, this is $\sim10^{-10}\,\mathrm{m\,s^{-2}}$, compatible with the
empirical BTFR calibration.
\end{corollary}

\section{Not MOND}
An effective $a_0$ appears, but the dynamics remain residual ordinary gravity with a
ballistic halo---not a MOND interpolating function.

\section{What is theorem vs assumption}
\begin{itemize}
\item Inherited: BTFR attractor structure from galactic modelling.
\item Literature: $\beta=2\pi/\kappa$, cosmological $\kappa=H_0$.
\item Theorem: $a_{\mathrm{IR}}=H_0c/2\pi$; final $a_0=H_0c/2\pi$ given the pins.
\item Physical inputs: far-halo IR pin; $a_0=a_{\mathrm{IR}}$.
\item Not claimed: MOND force law; separate microscopic $C$, $A$.
\end{itemize}

\end{document}
